\let\negmedspace\undefined
\let\negthickspace\undefined
\documentclass[journal]{IEEEtran}
\usepackage[a5paper, margin=10mm, onecolumn]{geometry}
%\usepackage{lmodern} % Ensure lmodern is loaded for pdflatex
\usepackage{tfrupee} % Include tfrupee package

\setlength{\headheight}{1cm} % Set the height of the header box
\setlength{\headsep}{0mm}     % Set the distance between the header box and the top of the text

\usepackage{gvv-book}
\usepackage{gvv}
\usepackage{cite}
\usepackage{amsmath,amssymb,amsfonts,amsthm}
\usepackage{algorithmic}
\usepackage{graphicx}
\usepackage{textcomp}
\usepackage{xcolor}
\usepackage{txfonts}
\usepackage{listings}
\usepackage{enumitem}
\usepackage{mathtools}
\usepackage{gensymb}
\usepackage{comment}
\usepackage[breaklinks=true]{hyperref}
\usepackage{tkz-euclide} 
\usepackage{listings}
% \usepackage{gvv}                                        
\def\inputGnumericTable{}                                 
\usepackage[latin1]{inputenc}                                
\usepackage{color}                                            
\usepackage{array}                                            
\usepackage{longtable}                                       
\usepackage{calc}                                             
\usepackage{multirow}                                         
\usepackage{hhline}                                           
\usepackage{ifthen}                                           
\usepackage{lscape}
\begin{document}
\bibliographystyle{IEEEtran}
\title{12.640}
\author{EE25BTECH11002 - Achat Parth Kalpesh }
{\let\newpage\relax\maketitle}
\renewcommand{\thefigure}{\theenumi}
\renewcommand{\thetable}{\theenumi}
\setlength{\intextsep}{10pt} % Space between text and floats
\numberwithin{equation}{enumi}
\numberwithin{figure}{enumi}
\renewcommand{\thetable}{\theenumi}
\parindent 0px



\textbf{Question:}\\
$\vec{P}$ and $\vec{R}$ are Jacobian matrices for two different geophysical inverse problems. If their generalized inverses are written as $\vec{P}^{-g} = \brak{\vec{P}^T \vec{P}}^{-1} \vec{P}^T$ and $\vec{R}^{-g} = \vec{R}^T\brak{\vec{R} \vec{R}^T}^{-1}$, then
\begin{enumerate}
    \item $\vec{P}$ represents an overdetermined problem and $\vec{R}$ represents an underdetermined problem
    \item $\vec{P}$ represents an underdetermined problem and $\vec{R}$ represents an overdetermined problem
    \item Both $\vec{P}$ and $\vec{R}$ represent overdetermined problems
    \item Both $\vec{P}$ and $\vec{R}$ represent underdetermined problems
\end{enumerate}

\textbf{Solution:}\\
Given are two Jacobian matrices $\vec{P}$ and $\vec{R}$ with generalized inverses:
\begin{align}
    \vec{P}^{-g} &= \brak{\vec{P}^T \vec{P}}^{-1} \vec{P}^T \\
    \vec{R}^{-g} &= \vec{R}^T \brak{\vec{R} \vec{R}^T}^{-1}
\end{align}
Let $\vec{P} \in \mathbb{R}^{m \times n}$ and $\vec{R} \in \mathbb{R}^{m \times n}$
\begin{align}
    \vec{P}^{-g} = \brak{\vec{P}^T \vec{P}}^{-1} \vec{P}^T\\
    \vec{P}^T \vec{P} \in \mathbb{R}^{n \times n}
\end{align}
For the inverse to exist
\begin{align}
   \det\brak{\vec{P}^T \vec{P}} \neq 0 \\
   \implies \text{rank}\brak{\vec{P}} = n \leq m
\end{align}
This implies $m \geq n$, i.e., more equations than unknowns. Hence, $\vec{P}$ is tall and the system is overdetermined
\begin{align}
    \vec{R}^{-g} = \vec{R}^T \brak{\vec{R} \vec{R}^T}^{-1}\\
    \vec{R} \vec{R}^T \in \mathbb{R}^{m \times m}
\end{align}
For the inverse to exist
\begin{align}
    \det\brak{\vec{R} \vec{R}^T} \neq 0 \\
    \implies \text{rank}\brak{\vec{R}} = m \leq n
\end{align}
This implies $n \geq m$, i.e., more unknowns than equations. Hence, $\vec{R}$ is wide and the system is underdetermined\\
$\vec{P}$ represents an overdetermined system and $\vec{R}$ represents an underdetermined system

\end{document}
